% =============================================================================
% Improved version of the Player B / Punishment Behavior subsection.
%
% Improvements vs the current results.tex L28-62:
%   - Fix "(\ref{fig:punishment})" → "Figure~\ref{...}, panel~A"
%   - Fix subject-verb agreement: "correlations ... are essentially zero"
%   - Fix the C4 footnote per the user's [Comment: the smallest is relevant]:
%     count subjects by their MIN non-zero per-cell punishment, not max.
%   - Drop the redundant in-line redefinition of (punish_nodel - punish_del)
%     in the predictors paragraph (defined two paragraphs above).
%   - Drop "confirming the bivariate result" sentence in the predictors
%     paragraph: same diff-in-diff, repeated.
%   - Update column references: full sample is cols 2 (and 3 with controls);
%     passer parallels are cols 5 (and 6 with controls).
%   - Drop the standalone "What predicts punishment?" paragraph and fold its
%     belief-check argument into the H2-finding paragraph (where the
%     regression table is already cited). The regression table is now
%     introduced once.
%   - Move the preregistered-exclusion footnote next to the actual robustness
%     check, so it doesn't appear two paragraphs before the check itself.
%   - Move the Result 2 box to immediately follow the conclusion paragraph.
%   - Tighten the conclusion / H1<->H2 bridge so the parallelism reads cleanly.
%   - Unify "scenario" and "cell" → "cell" throughout the body for consistency.
%   - Repair "punsihment" typo in figure caption.
%   - Drop stale [Comment: ...] markers that are now resolved or moved here.
%
% Note: the figure caption is rewritten to match the current 2-panel figure
% (Panel A actual, Panel B hypothetical, both dual y-axis).
% =============================================================================

\subsection*{Player~B -- Punishment behavior} \label{sec:result2}

We turn to Player~B's punishment behavior in the \textit{Punishment} condition, in which Player~B was given a real, payoff-relevant punishment choice. The default sample throughout is the full $N=80$ Punishment-condition sample; we report a robustness check on the $67$ subjects who passed the preregistered attention check on the punishment-elicitation screen below.

Punishment in all four (delegation, outcome) cells of the strategy method is significantly greater than zero (one-sample $t$-test $p<10^{-5}$ in every cell). Player~B punishes more for bad outcomes than for good outcomes---a within-subject mean difference of about $\pounds 0.21$, significant at the $1\%$ level under both paired $t$- and Wilcoxon signed-rank tests---an indication that punishment carries informational content rather than being random.\footnote{\citet{coffman_intermediation_2011} similarly documents non-trivial punishment in scenarios where the decision-maker's action could not have been improved upon. Positive punishment in the good-outcome cells therefore need not reflect noise: Player~Bs may hold ex-post views about whether the \textit{delegation choice itself} was appropriate, independent of the realised outcome, and can punish accordingly.} $45.0\%$ of Player~Bs do not punish in any cell.\footnote{Among the $44$ Player~Bs who do punish in at least one cell, the punishment levels are predominantly non-trivial: only $6$ ever choose a strictly positive cell punishment of at most $\pounds 0.10$, $14$ at most $\pounds 0.20$, and $19$ at most $\pounds 0.30$ (cumulative, smallest non-zero cell per subject). Conversely, $15$ punishers choose the maximum $\pounds 2$ in at least one cell. The costly punishment option is therefore being used as a substantive sanction rather than as tokenistic enforcement.}

Conditional on the outcome, Player~A is not significantly differently punished depending on the delegation decision (Figure~\ref{fig:punishment}, panel~A). When Player~B receives the high payoff, average punishment is $\pounds 0.329$ for delegated decisions and $\pounds 0.297$ for self-made decisions. When Player~B receives the low payoff, average punishment is $\pounds 0.482$ for delegated decisions and $\pounds 0.553$ for self-made decisions; this difference points in the direction of H2 but is not significant (paired $t$-test $p=0.131$; Wilcoxon signed-rank $p=0.544$). Equivalently, regressing the within-subject difference in punishment (self-made minus delegated) on a bad-outcome indicator yields $-\pounds 0.03$ in the good cell and $+\pounds 0.07$ in the bad cell, with the diff-in-diff between the two cells of $+\pounds 0.10$ borderline at $p=0.07$ (Table~\ref{tab:reg_punishment}, column~2; column~3 adds socio-demographic controls and Player~B's belief about Player~A's performance, leaving the diff-in-diff unchanged at $+\pounds 0.10$, $p=0.08$). The same picture arises for punishment frequencies (Figure~\ref{fig:punishment}, panel~A, right axis): Player~B neither punishes delegated decisions substantially less harshly, nor less frequently. One alternative driver of cell-level punishment we can also rule out is heterogeneity in Player~B's beliefs about Player~A's absolute performance: although Player~A and the algorithm are calibrated to perform identically and this calibration is common knowledge, Player~Bs hold heterogeneous absolute-level beliefs about Player~A's score, and those beliefs are uncorrelated with chosen punishment in every cell (Pearson $|r|\le 0.03$, all $p>0.8$).

The result is essentially identical when we restrict to the $67$ Player~Bs who passed the punishment-screen attention check (mean difference $\pounds 0.077$ vs $\pounds 0.071$; paired $t$-test $p=0.153$ vs $p=0.131$; Table~\ref{tab:reg_punishment}, columns~5--6).

In our experimental setting, Player~A therefore cannot reduce her expected punishment by delegating to the algorithm: punishment is determined almost entirely by the realised outcome and not by the delegation decision. This non-result also bears on \textit{Result~1}: a treatment difference in delegation rates motivated by responsibility avoidance would require Player~B to differentially punish delegated decisions, which we do not observe.

\medskip
\noindent\textbf{Result~2.} \textit{Player~A cannot avoid punishment by delegating the decision to the algorithm.}
\medskip

\begin{figure}[!htbp]
    \centering
    \caption{Player~B punishment by delegation and outcome cell}
    \includegraphics[width=\linewidth]{figures/punishment.png}
    \label{fig:punishment}
    \begin{minipage}{0.95\linewidth}
    \footnotesize
    \textit{Note:} Both panels report, for each of the four (delegation, outcome) cells of the strategy method, the average punishment chosen by Player~B (left $y$-axis, in pounds, error bars are standard errors of the mean) and the share of Player~Bs choosing zero punishment (right $y$-axis, in percent). In each panel, the full sample is shown at full opacity; the attention-check passers are overlaid at lower opacity. Panel~A reports the actual punishment in the \textit{Punishment} condition (full sample $n=80$, passers $n=67$). Panel~B reports the hypothetical punishment in the \textit{No-Punishment} condition (full sample $n=81$, passers $n=68$), where Player~Bs answered the same questions but could not impose punishment in practice.
    \end{minipage}
\end{figure}

\begin{table}[!htbp]
\centering
\caption{Determinants of Player~B's chosen punishment}
\label{tab:reg_punishment}
\begin{tabular}{@{\extracolsep{5pt}}lccc|ccc}
\\[-1.8ex]\hline
\hline \\[-1.8ex]
& \multicolumn{6}{c}{Dependent variable: \textit{Punishment}}\\[0.1cm]
& \multicolumn{3}{c}{\textit{Full sample}} & \multicolumn{3}{c}{\textit{Passers}} \\
& \textit{Level} (\pounds) & \multicolumn{2}{c}{\textit{Difference} (\pounds)} & \textit{Level} (\pounds) & \multicolumn{2}{c}{\textit{Difference} (\pounds)} \\
\\[-1.8ex] & (1) & (2) & (3) & (4) & (5) & (6) \\
\hline \\[-1.8ex]
 Delegated & $0.032$ &  &  & $0.001$ &  &  \\
  & $(0.052)$ &  &  & $(0.054)$ &  &  \\[0.1cm]
 Bad outcome & $0.256^{***}$ & $0.103^{*}$ & $0.103^{*}$ & $0.260^{***}$ & $0.078$ & $0.078$ \\
  & $(0.071)$ & $(0.057)$ & $(0.059)$ & $(0.080)$ & $(0.052)$ & $(0.053)$ \\[0.1cm]
 Delegated $\times$ Bad outcome & $-0.103^{*}$ &  &  & $-0.078$ &  &  \\
  & $(0.058)$ &  &  & $(0.052)$ &  &  \\[0.1cm]
 Player B belief about Player A perf. & $-0.020$ &  & $-0.006$ & $0.001$ &  & $-0.006$ \\
  & $(0.032)$ &  & $(0.029)$ & $(0.037)$ &  & $(0.038)$ \\[0.1cm]
 Age & $0.002$ &  & $-0.009^{**}$ & $-0.001$ &  & $-0.009^{**}$ \\
  & $(0.005)$ &  & $(0.003)$ & $(0.006)$ &  & $(0.004)$ \\[0.1cm]
 Female & $-0.100$ &  & $-0.012$ & $-0.064$ &  & $0.003$ \\
  & $(0.128)$ &  & $(0.109)$ & $(0.140)$ &  & $(0.132)$ \\[0.1cm]
 Socio-economic status & $0.083^{*}$ &  & $0.029$ & $0.072$ &  & $0.041$ \\
  & $(0.047)$ &  & $(0.024)$ & $(0.049)$ &  & $(0.028)$ \\[0.1cm]
 Went to uni & $-0.196$ &  & $0.051$ & $-0.098$ &  & $0.023$ \\
  & $(0.140)$ &  & $(0.070)$ & $(0.142)$ &  & $(0.075)$ \\[0.1cm]
 Technology score & $0.073$ &  & $-0.022$ & $0.030$ &  & $-0.012$ \\
  & $(0.053)$ &  & $(0.027)$ & $(0.059)$ &  & $(0.033)$ \\[0.1cm]
 Leadership position & $0.083$ &  & $0.143$ & $0.061$ &  & $0.174$ \\
  & $(0.142)$ &  & $(0.122)$ & $(0.164)$ &  & $(0.145)$ \\[0.1cm]
 Constant & $-0.077$ & $-0.032$ & $0.116$ & $0.020$ & $-0.001$ & $0.076$ \\
  & $(0.328)$ & $(0.051)$ & $(0.320)$ & $(0.374)$ & $(0.053)$ & $(0.397)$ \\[0.1cm]
\hline \\[-1.8ex]
 Observations & $320$ & $160$ & $160$ & $268$ & $134$ & $134$ \\
 Subjects & $80$ & $80$ & $80$ & $67$ & $67$ & $67$ \\
 Adj.\ $R^2$ & $0.077$ & $0.008$ & $0.038$ & $0.030$ & $0.001$ & $0.043$ \\
\hline
\hline \\[-1.8ex]
\multicolumn{7}{p{0.85\textwidth}}{\footnotesize \textit{Note:} Columns~(1) and~(4): OLS regression of Player~B's chosen punishment level on a delegation indicator, a bad-outcome indicator, their interaction, Player~B's belief about Player~A's performance, and a vector of socio-demographic controls; sample stacked over the four (delegation, outcome) cells of the strategy method (4 observations per Player~B). Columns~(2)/(5) and~(3)/(6): OLS regression of the within-subject delegation insulation (\textit{punish\_nodel} $-$ \textit{punish\_del}) on a bad-outcome indicator and, in Columns~(3)/(6), additional controls; sample stacked over the two outcome cells (2 observations per Player~B). In the difference specifications, the intercept is the average insulation in the good-outcome cell and the bad-outcome coefficient is the additional insulation in the bad-outcome cell. Columns~(1)--(3) use the full Punishment-condition sample. Columns~(4)--(6) restrict to Player~Bs who passed the attention check on the punishment-elicitation screen, as preregistered. Standard errors clustered at the Player~B level. Significance: $^{*}\,p<0.10$; $^{**}\,p<0.05$; $^{***}\,p<0.01$ (two-sided).}
\end{tabular}
\end{table}


% =============================================================================
% End of subsection.
% =============================================================================
